\section{Experiment}
\subsection{Evaluation}

\noindent\textbf{Evaluation Datasets.}
In the absence of an existing dataset containing both camera and object motion, we curated a diverse collection of 100 real-world videos from publicly available sources, carefully selected to cover a wide range of camera trajectories and dynamic object movements.


\noindent\textbf{Evaluation Metrics.}
Following prior work~\cite{wang2024motionctrl, li2025magicmotion, fu20243dtrajmaster, shuai2025free}, we evaluate performance across four key dimensions:
(1) Visual Quality. We use Fréchet Image Distance (FID)~\cite{maximilian2020fid} to evaluate visual fidelity and Fréchet Video Distance (FVD)~\cite{unterthiner2018towards} to assess temporal coherence.
(2) Text Alignment. CLIP Similarity (CLIPSIM)~\cite{wu2021godiva} measures the semantic consistency between each generated video and its corresponding text prompt.
(3) Camera Motion. Following CameraCtrl, we adopt CamTransErr and CamRotErr to quantify the translation and rotation deviations between the generated and reference camera trajectories.
(4) Object Motion. We use Box-IoU to evaluate the accuracy of object trajectories. For each generated video, we obtain predicted masks $M_{\text{gen}}$ by providing the ground-truth first-frame masks $M_{\text{gt}}(0)$ to SAM2~\cite{ravi2024sam}. Bounding boxes are then derived from each mask, and the mean Intersection-over-Union (IoU) between predicted and ground-truth boxes across all frames is reported as the final Box-IoU score.

\subsection{Comparisons with State-of-the-Art Methods}
We begin by evaluating SymphoMotion under independent camera control, comparing it against prior approaches~\cite{he2024cameractrl,yu2024viewcrafter,cao2025uni3c}. We then assess its capability for unified control, where both motions must be coordinated within a single framework. Across both settings, SymphoMotion consistently demonstrates superior controllability and fidelity.

\begin{figure}
    \centering
    \includegraphics[width=0.999\linewidth]{images/fig5-camera.pdf}
    \vspace{-7mm}
    \caption{\textbf{Independent camera motion control.}}
    \label{fig:camera-control}
\end{figure}

\noindent\textbf{Independent Control of Camera Motion.}
CameraCtrl~\cite{he2024cameractrl}, ViewCrafter~\cite{yu2024viewcrafter}, and Uni3C~\cite{cao2025uni3c} are selected for comparison, as all accept explicit camera specifications. As shown in Fig.~\ref{fig:camera-control}, when simulating controlled camera motion, CameraCtrl produces noticeable distortions, while ViewCrafter fails to preserve static objects during viewpoint changes. In contrast, Uni3C and SymphoMotion more faithfully reflect the intended camera behavior. The CamTransErr and CamRotErr metrics in Tab.~\ref{tab:performance_metrics} further indicate that SymphoMotion achieves camera-control accuracy comparable to these dedicated baselines.

\noindent\textbf{Simultaneous Control of Camera and Object Motions.}
\begin{figure}[t]
    \centering
    \includegraphics[width=0.97\linewidth]{images/fig6-camera-object.pdf}
    \caption{\textbf{Simultaneous control over camera and object motions.} MotionCtrl struggles to generate realistic object dynamics, causing objects to disappear from view, whereas SymphoMotion achieves high-quality simultaneous control.}
    \label{fig:simul_comparison}
\end{figure}
We evaluate joint control of camera and object motion by comparing our method with existing approaches. Since CameraCtrl, Uni3C, and ViewCrafter support only camera motion, we adopt MotionCtrl~\cite{wang2024motionctrl} as the primary baseline for object motion. As shown in Fig.~\ref{fig:simul_comparison}, videos generated by SymphoMotion align more faithfully with the specified control signals, producing coordinated camera movement and realistic object dynamics. In contrast, MotionCtrl exhibits limited controllability for both camera and object motion; its camera trajectories deviate from the prescribed path, and its object behavior remains inconsistent or implausible. Quantitatively, SymphoMotion attains higher Box-IoU scores in Tab.~\ref{tab:performance_metrics}, indicating more accurate object trajectory adherence. Furthermore, as shown in Tab.~\ref{tab:user_study}, SymphoMotion receives higher ratings in the user study, outperforming previous methods in visual and motion quality.

\begin{table}[t]
\centering
\caption{Quantitative comparison of our method SymphoMotion with CameraCtrl, ViewCrafter, Uni3C and MotionCtrl.}
\label{tab:performance_metrics}
\vspace{-3mm}
\scriptsize
\setlength{\tabcolsep}{0.9mm}
\begin{tabular}{l|c|c|c|c|c}
\specialrule{.1em}{.05em}{.05em}
\rowcolor{mygray!50}
Method & CameraCtrl & ViewCrafter & Uni3C & MotionCtrl & SymphoMotion \\
\hline\hline

FID $\downarrow$
& 196.84
& 303.83
& \underline{86.66}
& 182.15
& \textbf{70.47} \\

\rowcolor{mygray2!36}
FVD $\downarrow$
& 1019.49
& 1690.73
& \underline{404.21}
& 738.41
& \textbf{332.50} \\

CLIPSIM $\uparrow$
& 0.29
& 0.28
& \underline{0.31}
& 0.30
& \textbf{0.31} \\

\rowcolor{mygray2!36}
CamTransErr $\downarrow$
& 0.68
& 0.80
& \underline{0.44}
& 0.83
& \textbf{0.37} \\

CamRotErr $\downarrow$
& 0.12
& 0.21
& \underline{0.06}
& 0.23
& \textbf{0.05} \\

\rowcolor{mygray2!36}
Box-IoU $\uparrow$
& --
& --
& --
& \underline{31.42}
& \textbf{61.88} \\
\specialrule{.1em}{.05em}{.05em}
\end{tabular}
\vspace{-4mm}
\end{table}


\subsection{Ablation Studies}
\begin{table}[t]
\centering
\caption{User study on visual quality, text alignment, camera motion, and object motion (scores range from 1 to 5, higher is better).}
\label{tab:user_study}
\vspace{-3mm}
\footnotesize
\setlength{\tabcolsep}{0.96mm}
\resizebox{0.476\textwidth}{!}{
\begin{tabular}{l|c|c|c|c|c}
\specialrule{.1em}{.05em}{.05em} 
\rowcolor{mygray!50}
Method 
& CameraCtrl
& ViewCrafter
& Uni3C
& MotionCtrl 
& SymphoMotion \\
\hline\hline

Visual Quality 
& 3.43
& 4.03
& 4.24
& 3.53
& \textbf{4.87} \\

\rowcolor{mygray2!36}
Text Alignment
& 3.37
& 3.57
& 3.68
& 3.14
& \textbf{4.02} \\

Camera Motion
& 3.13
& 3.47
& 3.97
& 3.18
& \textbf{4.36} \\

\rowcolor{mygray2!36}
Object Motion
& --
& --
& --
& 2.87
& \textbf{4.58} \\

\specialrule{.1em}{.05em}{.05em} 
\end{tabular}
}
\vspace{-1mm}
\end{table}



\begin{table}[t]
\centering
\caption{Quantitative results in ablation study.}
\label{tab:ablation_quantitative}
\vspace{-3mm}
\footnotesize
\setlength{\tabcolsep}{0.6mm}
\resizebox{0.47\textwidth}{!}{
\begin{tabular}{l|c|c|c|c}
\specialrule{.1em}{.05em}{.05em} 
\rowcolor{mygray!50}Metrics & FVD$\downarrow$ & CamTransErr$\downarrow$ & CamRotErr$\downarrow$ & Box-IoU$\uparrow$ \\
\hline\hline

\rowcolor{mygray2!36}
\textit{w/o} $c_{pcd}$ 
& 330.64 & 0.46 & 0.07 & 56.74 \\

\textit{w/o} 2D boxes  
& 337.14 & \textbf{0.36} & 0.06 & 54.32 \\

\rowcolor{mygray2!36}
\textit{w/o} 3D trajectory 
& 343.80 & \textbf{0.36} & 0.06 & 52.16 \\

SymphoMotion 
& \textbf{332.50} & 0.37 & \textbf{0.05} & \textbf{61.88} \\

\specialrule{.1em}{.05em}{.05em} 
\end{tabular}
}
\vspace{-1mm}
\end{table}


\begin{figure}
    \centering
    \includegraphics[width=0.99\linewidth]{images/fig7-ablation.pdf}
    \vspace{-2mm}
    \caption{\textbf{Results of different settings in the ablation study.}}
    \label{fig:ablation}
    \vspace{-1mm}
\end{figure}


\noindent\textbf{Effect of 3D Geometric Priors in Camera Trajectory Control.}
As shown in Tab.~\ref{tab:ablation_quantitative} and the first row of Fig.~\ref{fig:ablation}, incorporating geometry-aware renderings yields more accurate camera trajectories due to additional 3D structural priors conveyed by these renderings, which provide richer spatial understanding and informative viewpoint cues. By supplying complementary spatial cues, the geometry-aware renderings help the model better disambiguate camera motion, producing videos that closely follow the target trajectory and exhibit higher geometric consistency.

\vspace{1mm}
\noindent \textbf{Effect of 2D Visual Guidance in Object Dynamics Control.}
2D visual guidance provides an anchor for each object’s projected motion in the image plane. We find that incorporating these cues into the rendered frames enhances training stability and strengthens the model’s ability to reason about object motion. As shown in the second row of Fig.~\ref{fig:ablation}, introducing 2D visual guidance enables the model to better interpret object trajectories, producing videos with more coherent and physically plausible motion.

\vspace{1mm}
\noindent \textbf{Effect of 3D Trajectory Conditioning in Object Dynamics Control.}
3D trajectories provide explicit guidance for modeling object motion in space. As shown in Tab.~\ref{tab:ablation_quantitative}, incorporating object-level 3D trajectories leads to consistently better performance on object-motion metrics. As illustrated in the third row of Fig.~\ref{fig:ablation}, when the model is conditioned on a precise 3D spatial trajectory for the duck, it produces more faithful object motion, such as realistic head turns, indicating a stronger ability to follow complex spatial trajectories.

\section{Conclusion}
This paper presents SymphoMotion, a unified framework for jointly controlling camera motion and object dynamics in video generation. The framework builds on two key mechanisms effectively enabling unified motion control. Camera Trajectory Control leverages pose conditioning with 3D structural information for precise and stable viewpoint manipulation. Object Dynamics Control combines 2D spatial signals with 3D motion representations, allowing objects to follow intended trajectories with accurate motion. Due to the lack of high-quality data, we construct RealCOD-25K to support motion control research. Extensive experiments show that SymphoMotion synthesizes realistic and temporally coherent videos that consistently follow user-specified camera motions and object trajectories.